\chapter{Methodology}

\section{System Overview}

The ANON4REID evaluation framework consists of three stages: (1)~anonymization of the Market-1501 dataset at four levels of increasing privacy strength, (2)~utility evaluation by measuring \acrshort{reid} performance of a pre-trained \acrshort{osnet} model on each anonymized variant, and (3)~privacy evaluation by measuring the identity classification accuracy of an adversarial MobileNetV2 attacker on the anonymized query images. Figure~\ref{fig:pipeline} illustrates the overall pipeline.

\begin{figure}[H]
    \centering
    \begin{tikzpicture}[
        node distance=0.6cm and 1.2cm,
        box/.style={rectangle, draw=cfDark, fill=cfLightGray, thick, minimum height=0.9cm, minimum width=2.4cm, align=center, font=\small},
        anonbox/.style={rectangle, draw=cfAmber, fill=cfCream, thick, minimum height=0.9cm, minimum width=2.8cm, align=center, font=\small},
        evalbox/.style={rectangle, draw=cfTeal, fill=cfTeal!10, thick, minimum height=0.9cm, minimum width=2.4cm, align=center, font=\small},
        resultbox/.style={rectangle, draw=cfOrange, fill=cfOrange!10, thick, minimum height=0.9cm, minimum width=2.4cm, align=center, font=\small},
        arr/.style={-{Stealth[length=3mm]}, thick, cfDark},
    ]
        % Input
        \node[box] (input) {Market-1501\\(36,036 images)};

        % Anonymization levels
        \node[anonbox, right=1.5cm of input, yshift=1.8cm]  (l1) {L1: Blur};
        \node[anonbox, right=1.5cm of input, yshift=0.6cm]  (l2) {L2: SD Inpaint};
        \node[anonbox, right=1.5cm of input, yshift=-0.6cm] (l3) {L3: RAD};
        \node[anonbox, right=1.5cm of input, yshift=-1.8cm] (l4) {L4: Edge};

        % Evaluation
        \node[evalbox, right=1.8cm of l2, yshift=0.6cm]  (utility) {Utility Eval\\(OSNet)};
        \node[evalbox, right=1.8cm of l3, yshift=-0.6cm]   (privacy) {Privacy Eval\\(MobileNetV2)};

        % Output
        \node[resultbox, right=1.2cm of utility, yshift=-1.2cm] (result) {Trade-off\\Analysis};

        % Arrows: input -> levels
        \draw[arr] (input.east) -- ++(0.4,0) |- (l1.west);
        \draw[arr] (input.east) -- ++(0.4,0) |- (l2.west);
        \draw[arr] (input.east) -- ++(0.4,0) |- (l3.west);
        \draw[arr] (input.east) -- ++(0.4,0) |- (l4.west);

        % Arrows: levels -> evaluations
        \draw[arr] (l1.east) -- ++(0.6,0) |- (utility.west);
        \draw[arr] (l2.east) -- ++(0.6,0) |- (utility.west);
        \draw[arr] (l3.east) -- ++(0.6,0) |- (privacy.west);
        \draw[arr] (l4.east) -- ++(0.6,0) |- (privacy.west);

        % Arrows: evaluations -> result
        \draw[arr] (utility.east) -- ++(0.3,0) |- (result.west);
        \draw[arr] (privacy.east) -- ++(0.3,0) |- (result.west);

    \end{tikzpicture}
    \caption{ANON4REID evaluation pipeline. Each anonymized dataset variant is evaluated for both utility (OSNet \acrshort{reid} performance) and privacy (MobileNetV2 attacker accuracy).}
    \label{fig:pipeline}
\end{figure}

The anonymization levels are ordered by intended privacy strength:

\begin{table}[H]
    \centering
    \begin{tabular}{@{}clll@{}}
        \toprule
        \rowcolor{cfTeal} \thdr{Level} & \thdr{Name} & \thdr{Technique} & \thdr{Scope} \\
        \midrule
        L1 & Gaussian Blur & MediaPipe face detection + blur & Face region \\
        L2 & SD Inpaint & Stable Diffusion inpainting & Central torso \\
        L3 & RAD & ControlNet OpenPose + SD v1.5 & Full body replacement \\
        L4 & Edge Silhouette & Canny edge detection & Entire image \\
        \bottomrule
    \end{tabular}
    \caption{Summary of the four anonymization levels.}
    \label{tab:anon_levels}
\end{table}


\section{Anonymization Pipeline}

\subsection{Level 1: Gaussian Blur with Face Detection}

Level~1 applies Gaussian blur (kernel size $15 \times 15$, $\sigma = 30$) to the face region of each pedestrian image. Face localization uses the MediaPipe Face Detection model \cite{lugaresi2019mediapipe}, specifically the short-range model variant with a minimum detection confidence threshold of 0.3.

Because Market-1501 images are only 64$\times$128 pixels, face detection frequently fails due to insufficient resolution. A spatial fallback is therefore implemented: when MediaPipe returns no detections, the blur is applied to a heuristically defined head region covering the top 28\% of the image height and the central 70\% of the image width (from 15\% to 85\%). This fallback is based on the observation that Market-1501 bounding boxes are consistently cropped around the full body, placing the head in the upper portion of the frame.

\begin{equation}
    \text{ROI}_\text{fallback} = \text{image}[0 : \lfloor 0.28 \cdot H \rfloor,\; \lfloor 0.15 \cdot W \rfloor : \lfloor 0.85 \cdot W \rfloor]
    \label{eq:blur_fallback}
\end{equation}

An alternative design would apply the upscale--process--downscale pipeline (Section~\ref{sec:resolution}) to Level~1, giving MediaPipe higher-resolution input for face detection. This approach was considered and rejected for several reasons:

\begin{enumerate}[noitemsep]
    \item Lanczos upscaling interpolates between existing pixels without adding genuine facial detail. The upscaled ``face'' remains a smoothed approximation of the 64$\times$128 original, so MediaPipe may still fail on synthetic features that lack real structure.
    \item The spatial fallback already produces reasonable results for Market-1501 specifically, because the dataset's bounding boxes are consistently cropped around full-body pedestrians with heads centered in the upper portion of the frame.
    \item Adding upscaling would introduce per-image resize overhead to what is intended to be the lightweight, CPU-bound anonymization level. Part of Level~1's value in the benchmark is its minimal computational cost.
    \item More accurate face detection could produce tighter blur regions than the current fallback crop. A smaller blurred area might actually \textit{increase} facial information leakage at the region boundaries compared to the more conservative heuristic crop.
    \item The core finding is unaffected: the attacker exploits body-level features (clothing, shape, accessories) that face blurring leaves untouched regardless of localization precision. The 88.66\% identity leakage would persist even with perfect face detection.
\end{enumerate}


\subsection{Resolution Handling}
\label{sec:resolution}

Market-1501 images are natively 64$\times$128 pixels, far below the minimum resolution required by generative models. To address this, Levels~2 and~3 apply an upscale--process--downscale pipeline: each image is resized from 64$\times$128 to 256$\times$512 using Lanczos interpolation, the generative model operates on the upscaled image, and the output is resized back to 64$\times$128. This ensures coherent generative output while maintaining data compatibility with the original dataset structure.


\subsection{Level 2: Stable Diffusion Inpainting}

Level~2 uses the \texttt{sd-legacy/stable-diffusion-inpainting} model \cite{rombach2022ldm} to replace the central torso region with AI-generated content. A geometric binary mask covers the central 60\% width and 80\% height of the image (from 20\% to 80\% width, 10\% to 90\% height), targeting the body region that carries the most discriminative appearance features.

The inpainting model is conditioned on two text prompts:
\begin{itemize}[noitemsep]
    \item \textbf{Prompt:} \texttt{"empty background, no person, floor, wall"}
    \item \textbf{Negative prompt:} \texttt{"person, human, body, blurry, distorted"}
\end{itemize}

Inference uses 20 denoising steps with a guidance scale of 7.5. Each image is upscaled before processing and downscaled back afterward (Section~\ref{sec:resolution}).

\subsection{Level 3: Realistic Anonymization by Diffusion (RAD)}

Level~3 replaces the original person entirely with a synthetically generated individual who preserves the original body pose but shares no visual identity features. The pipeline consists of three components:

\begin{enumerate}[noitemsep]
    \item \textbf{Pose Extraction:} OpenPose \cite{cao2017openpose} extracts a skeleton representation of the person's body joints from the upscaled image, providing a privacy-preserving conditioning signal.
    \item \textbf{Conditioned Generation:} Stable Diffusion v1.5 \cite{rombach2022ldm} with ControlNet OpenPose conditioning \cite{zhang2023controlnet} generates a new person image guided by the extracted skeleton.
    \item \textbf{Inference Acceleration:} LCM-LoRA \cite{luo2023lcm} reduces the required denoising steps from 50 to 8, making batch processing of 36,036 images practical.
\end{enumerate}

The generation uses a guidance scale of 1.5 and a ControlNet conditioning scale of 0.8, with the prompt \texttt{"a person, different appearance, full body, high quality, realistic"}.

\subsection{Level 4: Canny Edge Detection}

Level~4 applies Canny edge detection \cite{canny1986edge} with thresholds $(T_\text{low}, T_\text{high}) = (100, 200)$ to the entire image, converting it to a binary edge map. The single-channel output is replicated to three channels (RGB) for compatibility with downstream models. This technique strips all color, texture, and fine-grained appearance information, retaining only structural contours.


\section{Evaluation Framework}

\textbf{Utility} is measured by the \acrshort{reid} performance of a pre-trained \acrshort{osnet} (\texttt{osnet\_x1\_0}) model evaluated on each anonymized dataset variant. The evaluation follows the standard closed-set \acrshort{reid} protocol: features are extracted from all query and gallery images as 512-dimensional embeddings, a Euclidean distance matrix is computed between all query--gallery pairs, and k-reciprocal encoding re-ranking \cite{zhong2017reranking} is applied to refine the distance matrix before computing \acrshort{cmc} curves and \acrshort{map} using the \texttt{torchreid} evaluation function with a maximum rank of 50. Each anonymized dataset variant is dynamically registered as a \texttt{torchreid} dataset using Python's \texttt{type()} to create subclasses of \texttt{Market1501} with modified directory paths, ensuring full compatibility with the standard evaluation pipeline without requiring custom dataloaders.

\textbf{Privacy} is measured by the Top-1 accuracy of an adversarial MobileNetV2 \cite{sandler2018mobilenetv2} classifier on anonymized query images. The attacker is pre-trained on ImageNet \cite{deng2009imagenet} and fine-tuned on original (non-anonymized) Market-1501 images to classify person identity. During inference, each anonymized query image is resized to 256$\times$128 pixels and normalized using the ImageNet channel statistics ($\mu = [0.485, 0.456, 0.406]$, $\sigma = [0.229, 0.224, 0.225]$) before being passed through the classifier. The privacy metric is:

\begin{equation}
    \text{Privacy Leakage} = \frac{\text{Correctly classified query images}}{\text{Total query images}} \times 100\%
    \label{eq:privacy_leak}
\end{equation}

A high leakage value indicates that the anonymization technique fails to protect identity information; a value near 0\% indicates strong privacy protection. For context, random guessing across 750 identity classes would yield approximately $\frac{1}{750} \times 100\% \approx 0.13\%$ accuracy, providing a theoretical floor for the privacy metric.

The attacker is trained on the \texttt{bounding\_box\_test} split (750 identities, 19,732 images) rather than \texttt{bounding\_box\_train}. This is because Market-1501's training and test/query sets have \textbf{completely disjoint identity sets} (751 vs.\ 750 IDs). Training on the training split would yield 0\% accuracy on query images regardless of anonymization, since none of the training identities appear in the query set. By training on the test split, which shares the same 750 identities as the query set, the attacker can meaningfully classify query images, and its accuracy drop on anonymized queries isolates the effect of the anonymization technique from the identity partition boundary.

\textbf{Trade-off.} The central analysis jointly examines utility degradation and privacy gain across anonymization levels. An effective anonymization technique should minimize utility loss while maximizing privacy protection. We define the \textit{Effectiveness Index} as:

\begin{equation}
    \text{Effectiveness Index} = \frac{\text{Privacy Protection (\%)}}{100 - \text{Utility (mAP \%)}}
    \label{eq:effectiveness}
\end{equation}

This ratio captures how much privacy is gained per unit of utility sacrificed. A higher value indicates a more efficient trade-off, while a low value indicates that privacy was gained at disproportionate utility cost.
